\documentclass[11pt,a4paper]{article}
\usepackage[margin=2.7cm]{geometry}
\usepackage[T1]{fontenc}
\usepackage[utf8]{inputenc}
\usepackage{lmodern}
\usepackage{amsmath,amssymb,amsthm,mathtools}
\usepackage{booktabs,array}
\usepackage{microtype}
\usepackage[hidelinks]{hyperref}
\usepackage{xcolor}
\usepackage{enumitem}
\usepackage{bm}

\newtheorem{theorem}{Theorem}[section]
\newtheorem{proposition}[theorem]{Proposition}
\newtheorem{lemma}[theorem]{Lemma}
\newtheorem{corollary}[theorem]{Corollary}
\theoremstyle{definition}
\newtheorem{definition}[theorem]{Definition}
\theoremstyle{remark}
\newtheorem{remark}[theorem]{Remark}

\newcommand{\PP}{\mathbb P}
\newcommand{\NN}{\mathbb N}
\newcommand{\OO}{\mathbb O}
\newcommand{\HH}{\mathbb H}
\newcommand{\RR}{\mathbb R}
\newcommand{\Sym}{\operatorname{Sym}}
\newcommand{\Aut}{\operatorname{Aut}}
\newcommand{\Span}{\operatorname{span}}
\newcommand{\divisor}{\operatorname{div}}

\title{\textbf{Arithmetic Occupation Spectra and the First Octonionic Nonassociativity Threshold}}
\author{EABC Research Note}
\date{September 2026}

\begin{document}
\maketitle

\begin{abstract}
We introduce a canonical occupation-space representation of natural numbers based solely on unique prime factorization. For
\[
N=\prod_{p\in\PP}p^{\nu_p},
\]
the associated factor word is
\[
W(N)=\sum_p\nu_p e_p,
\]
with arithmetic occupation numbers
\[
\Omega(N)=\sum_p\nu_p,
\qquad
\omega(N)=\#\{p:\nu_p>0\}.
\]
The logarithmic diagonal operator has exact spectral value
\[
E_\times(N)=\sum_p\nu_p\log p=\log N.
\]
Prime powers are therefore repeated occupations of one atomic prime mode, while products of distinct primes are genuine multimode states. We prove that in the sector $\Omega=3$ the three occupation types $p^3$, $p^2q$, and $pqr$ form an exact hierarchy: by alternativity and Artin's theorem, genuine three-generator octonionic nonassociativity is impossible for the first two types and becomes possible for the first time in the squarefree type $pqr$. We then separate the Euclidean tetrahedron volume from the $G_2$ calibration and the octonionic associator defect,
\[
V_{\rm tet}=\frac16\sqrt{\det G},\qquad
\varphi(u,v,w)=\langle uv,w\rangle,\qquad
\mathfrak t=\frac14\|[u,v,w]\|^2=\det G-\varphi^2.
\]
No dynamical prime-orbit realization or physical energy interpretation is assumed.
\end{abstract}

\section{Factor-word normal form}
Let
\[
\mathfrak M=\bigoplus_{p\in\PP}\NN_0 e_p
\]
be the free commutative monoid on the prime labels $e_p$. By unique prime factorization every $N\ge1$ has a unique representation
\[
N=\prod_{p\in\PP}p^{\nu_p}
\]
with finitely many nonzero $\nu_p$, and hence a canonical factor word
\begin{equation}
W(N)=\sum_p\nu_p e_p.
\end{equation}
Composite partial products create no new generators. Thus
\[
W(pq)=e_p+e_q,\qquad
W(p^2q)=2e_p+e_q,\qquad
W(pqr)=e_p+e_q+e_r.
\]

\begin{definition}[Arithmetic occupation numbers]
Define
\[
\Omega(N)=\sum_p\nu_p,
\qquad
\omega(N)=\#\{p:\nu_p>0\}.
\]
Thus $\Omega$ counts prime factors with multiplicity, while $\omega$ counts distinct occupied prime modes.
\end{definition}

\section{Logarithmic occupation spectrum}
Consider the Hilbert basis $\{|N\rangle:N\in\NN\}$ indexed by factor words. Define the diagonal operator
\begin{equation}
H_\times|N\rangle
=\left(\sum_p\nu_p\log p\right)|N\rangle.
\end{equation}

\begin{proposition}[Exact arithmetic energy]
For every $N\in\NN$,
\[
E_\times(N)=\log N.
\]
Moreover,
\[
E_\times(MN)=E_\times(M)+E_\times(N)
\]
for all $M,N\in\NN$.
\end{proposition}

\begin{proof}
Immediate from
\[
\log\left(\prod_pp^{\nu_p}\right)=\sum_p\nu_p\log p
\]
and $\log(MN)=\log M+\log N$.
\end{proof}

The atomic prime-mode energies are therefore
\[
\varepsilon_p=\log p,
\]
and a general occupation state has the free additive spectral value
\[
E_\times(N)=\sum_p\nu_p\varepsilon_p.
\]
This is an arithmetic spectral statement, not yet a physical energy claim.

\section{Prime, semiprime, and triple occupation sectors}
For $\Omega=1$ one has only prime states
\[
p:\quad W(p)=e_p,\qquad E_\times(p)=\log p.
\]
For $\Omega=2$ there are exactly two occupation types:
\[
p^2:\quad W=2e_p,\quad \omega=1,\quad E_\times=2\log p,
\]
\[
pq\ (p\ne q):\quad W=e_p+e_q,\quad \omega=2,\quad E_\times=\log p+\log q.
\]
Thus $p^2$ is repeated occupation of one mode, whereas $pq$ is a genuine two-mode state.

For $\Omega=3$ unique factorization yields exactly the partition types $3$, $2+1$, and $1+1+1$:
\begin{center}
\begin{tabular}{@{}llll@{}}
\toprule
Arithmetic type & Factor word & $\omega(N)$ & $E_\times(N)$ \\
\midrule
$p^3$ & $3e_p$ & $1$ & $3\log p$ \\
$p^2q$, $p\ne q$ & $2e_p+e_q$ & $2$ & $2\log p+\log q$ \\
$pqr$, pairwise distinct & $e_p+e_q+e_r$ & $3$ & $\log p+\log q+\log r$ \\
\bottomrule
\end{tabular}
\end{center}

If one uses bosonic occupation notation, these correspond formally to
\[
|3_p\rangle,\qquad |2_p,1_q\rangle,\qquad |1_p,1_q,1_r\rangle
\]
in $\Sym^3\mathcal H_1$. This symmetric occupation structure must not be identified with exterior powers appearing in Fredholm determinant expansions.

\section{Octonionic realization of occupied modes}
Assume that each distinct occupied prime mode is assigned a fixed imaginary octonionic direction
\[
p\longmapsto u_p\in\operatorname{Im}\OO.
\]
No canonicity of this assignment is assumed.

For $x,y,z\in\OO$, define the associator
\[
[x,y,z]=(xy)z-x(yz).
\]
The octonions are alternative, so
\[
[x,x,y]=[y,x,x]=0.
\]
Moreover, Artin's theorem implies that every subalgebra generated by two elements of an alternative algebra is associative.

\begin{theorem}[Occupation--associativity threshold]\label{thm:threshold}
Let $N$ satisfy $\Omega(N)=3$, and assign an octonionic generator to every distinct occupied prime mode. Then:
\begin{enumerate}[label=(\roman*)]
\item for $N=p^3$, every local three-word associator vanishes;
\item for $N=p^2q$ with $p\ne q$, every associator built from the two occupied generators $u_p,u_q$ vanishes;
\item for $N=pqr$ with $p,q,r$ pairwise distinct, a nonzero associator $[u_p,u_q,u_r]$ becomes possible.
\end{enumerate}
Consequently,
\[
\omega(N)\le2
\quad\Longrightarrow\quad
\text{genuine three-generator nonassociativity is impossible},
\]
whereas $\omega(N)=3$ is the first occupation threshold at which it may occur.
\end{theorem}

\begin{proof}
For $p^3$, all three entries are the same generator, hence alternativity gives $[u_p,u_p,u_p]=0$. For $p^2q$, at most the two generators $u_p,u_q$ occur. By Artin's theorem the subalgebra generated by them is associative, so every three-word associator formed from these generators vanishes. For $pqr$, three distinct generators occur and Artin's two-generator theorem no longer forces associativity. Thus nonzero associators are possible, though not necessary.
\end{proof}

\begin{remark}
The squarefree condition does not imply nonassociativity. Three distinct octonionic generators can still lie in a common associative quaternionic subalgebra. The theorem identifies the first \emph{possible} nonassociative occupation type, not a sufficient arithmetical criterion.
\end{remark}

\section{$G_2$ geometry of the three-mode configuration}
Let $u,v,w\in\operatorname{Im}\OO\cong\RR^7$. Write
\[
G=(\langle u_i,u_j\rangle)_{i,j=1}^3
\]
for their Gram matrix. The Euclidean volume of the simplex $\{0,u,v,w\}$ is
\begin{equation}
V_{\rm tet}(u,v,w)=\frac16\sqrt{\det G}.
\end{equation}
This is a purely metric invariant and must not be confused with the associator defect.

The standard associative $G_2$ three-form is, up to orientation convention,
\begin{equation}
\varphi(u,v,w)=\langle uv,w\rangle.
\end{equation}
Define the nonassociativity defect
\begin{equation}
\mathfrak t(u,v,w)=\frac14\|[u,v,w]\|^2.
\end{equation}
For imaginary octonions one has the identity
\begin{equation}\label{eq:g2identity}
\boxed{\mathfrak t=\det G-\varphi^2.}
\end{equation}
Equivalently,
\begin{equation}
\boxed{\mathfrak t=36V_{\rm tet}^2-\varphi^2.}
\end{equation}
Thus $V_{\rm tet}$, $\varphi$, and $\mathfrak t$ are related observables, not three independent invariants.

For an orthonormal triple, $G=I_3$ and hence
\[
V_{\rm tet}=\frac16,
\qquad
\mathfrak t=1-\varphi^2,
\qquad
0\le\mathfrak t\le1.
\]
General orthonormal $G_2$ triples need not be binary: intermediate values of $|\varphi|$ and $\mathfrak t$ are possible.

\section{Discrete Fano basis sector}
Let $e_1,\dots,e_7$ be the seven imaginary basis units in a fixed octonion multiplication table. There are
\[
\binom73=35
\]
unordered triples of distinct basis directions. Exactly seven are Fano-line triples and twenty-eight are non-Fano triples.

Every basis triple is orthonormal, hence all thirty-five associated Euclidean tetrahedra satisfy
\[
V_{\rm tet}=\frac16.
\]
The $G_2$ multiplication nevertheless separates them:
\begin{center}
\begin{tabular}{@{}lcccc@{}}
\toprule
Type & Number & $\det G$ & $|\varphi|$ & $\mathfrak t$ \\
\midrule
Fano line & $7$ & $1$ & $1$ & $0$ \\
Non-Fano basis triple & $28$ & $1$ & $0$ & $1$ \\
\bottomrule
\end{tabular}
\end{center}
Thus the thirty-five basis tetrahedra are Euclidean-volume equivalent but octonionically inequivalent. The distinction lies in the calibrated multiplication structure, not in ordinary metric volume.

\section{Configuration invariants versus arithmetic invariants}
For a fixed triple $(u_p,u_q,u_r)$ the quantities
\[
G,\qquad V_{\rm tet},\qquad\varphi,\qquad\mathfrak t
\]
are mathematically well-defined. They are therefore configuration invariants.

They are not yet invariants of the bare prime triple $(p,q,r)$. To obtain such an arithmetic invariant one would need a canonical or quotient-stable assignment
\[
p\longmapsto u_p
\]
or an invariant averaging procedure over a mathematically specified representation fiber. A choice of one four-square representation
\[
p=a^2+b^2+c^2+d^2
\]
and an arbitrary embedding of those coordinates into $\operatorname{Im}\OO$ does not supply such canonicity.

In particular, quotienting a single unit vector by the full group $G_2=\Aut(\OO)$ does not solve the problem: $G_2$ acts transitively on the unit sphere $S^6\subset\operatorname{Im}\OO$. Relevant $G_2$ invariants arise only for multi-vector configurations.

\section{Prime powers, multimode states, and possible orbit realizations}
The exact arithmetic identity
\[
E_\times(p^k)=k\log p
\]
exhibits prime powers as the occupation ladder
\[
|1_p\rangle,\ |2_p\rangle,\ |3_p\rangle,\dots
\]
along one prime mode. Mixed products activate several modes.

The formula resembles the repetition law
\[
\ell(\gamma^k)=k\ell(\gamma)
\]
for periodic orbits. However, the existence of a dynamics with
\[
p\mapsto\gamma_p,\qquad \ell(\gamma_p)=\log p
\]
is a separate realization problem and is not assumed here. In particular, the arithmetic equality $E_\times(p)=\log p$ must not be identified with a physically or dynamically realized primitive period without proof.

Likewise, bosonic occupation sectors
\[
\Sym^k\mathcal H_1
\]
must not be conflated with Fredholm exterior powers
\[
\Lambda^k\mathcal H.
\]
Any bridge between these constructions requires an independent theorem.

\section{Epistemic status and open gates}
The following statements are established within the algebraic framework:
\begin{itemize}
\item the canonical factor word $W(N)=\sum_p\nu_pe_p$;
\item the arithmetic occupation data $\Omega(N)$ and $\omega(N)$;
\item the exact logarithmic spectrum $E_\times(N)=\log N$;
\item the complete $\Omega=3$ decomposition into $p^3$, $p^2q$, and $pqr$;
\item Theorem~\ref{thm:threshold};
\item for a fixed octonionic triple, the geometric quantities $G$, $V_{\rm tet}$, $\varphi$, and $\mathfrak t$ and identity~\eqref{eq:g2identity};
\item the discrete $35=7+28$ Fano-basis classification.
\end{itemize}

The following remain open realization problems:
\begin{enumerate}
\item a canonical arithmetic assignment $p\mapsto u_p$;
\item representation-independence of the octonionic configuration invariants;
\item a periodic-orbit realization with primitive periods $\log p$;
\item a physical interpretation of $E_\times$ as measured energy;
\item any identification between factor-word Fock occupation and Fredholm exterior-power combinatorics.
\end{enumerate}

\section{Conclusion}
Unique prime factorization determines a canonical occupation vector, and logarithmic evaluation converts multiplicative arithmetic into an exactly additive spectrum:
\[
\boxed{
N=\prod_pp^{\nu_p}
\quad\longmapsto\quad
W(N)=\sum_p\nu_pe_p,
\qquad
E_\times(N)=\sum_p\nu_p\log p=\log N.
}
\]
Prime powers encode repeated occupation of one atomic mode, whereas mixed products encode genuine mode multiplicity. At occupation number three, alternativity and Artin's theorem force the $p^3$ and $p^2q$ sectors to remain associative and identify the squarefree $pqr$ sector as the first one in which a genuine octonionic associator can occur.

The resulting structural principle is
\[
\boxed{
\text{arithmetic mode complexity}
\;\Longrightarrow\;
\text{first possible algebraic nonassociativity threshold}.
}
\]
This conclusion is entirely algebraic. The remaining task is the arithmetic anchoring of the octonionic configuration, not the internal $G_2$ geometry itself.

\end{document}
